\chapter{用語と公式}\label{章：用語と計算の参照}\label{章：前提知識と計算の基本}\label{章：計算の基本}
\begin{bookcolumns}
本文で使う言葉と公式をまとめる.必要な項目だけ確かめ,参照先から元の議論へ戻れる.
和と式変形は基本型で,対数・三角関数・極限・複素数はそれぞれを使う解法で必要になる.

\section{用語一覧}\label{節：用語の対応}
本文で見た操作や性質には,一般に使われる名前がある.
ここには,その意味を短く記した.詳しい定義・条件・計算は,それぞれの参照先にある.
上から順に覚えるための一覧ではなく,必要な言葉から本文へ戻るための対応表である.
\begin{description}[style=nextline,leftmargin=0pt,labelsep=0pt,itemsep=5pt]
\item[\hypertarget{term:recurrence}{漸化式}] 項どうしの関係によって数列を定める式. \booknavref{節：数列と漸化式}{本文の説明へ.}
\item[\hypertarget{term:sequence}{数列}] 数を順に並べたもの.添字で各項の位置を表す. \booknavref{節：数列と漸化式}{本文の説明へ.}
\item[\hypertarget{term:initial}{初期条件}] 数列を作り始めるために与える値. \booknavref{節：数列と漸化式}{本文の説明へ.}
\item[\hypertarget{term:general}{一般項}] 番号からその位置の項を直接表す式. \booknavref{節：数列と漸化式}{本文の説明へ.}
\item[\hypertarget{term:sigma}{シグマ}] 和を,添字の範囲とともに書く記号. \booknavref{節：シグマの計算}{本文の説明へ.}
\item[\hypertarget{term:product}{総積}] 指定した範囲の項をすべて掛けたもの. \booknavref{節：階比型}{本文の説明へ.}
\item[\hypertarget{term:partial}{部分分数分解}] 一つの分数を,分母がより簡単な分数の和へ直す操作. \booknavref{節：部分分数分解}{本文の説明へ.}
\item[\hypertarget{term:square}{平方完成}] 同じ数を足して引くことで,二乗の形を作る操作. \booknavref{節：+a-aをする}{本文の説明へ.}
\item[\hypertarget{term:arithmetic}{等差数列}] 次の項へ進むとき,同じ数を足す数列. \booknavref{節：等差型}{本文の説明へ.}
\item[\hypertarget{term:geometric}{等比数列}] 次の項へ進むとき,同じ数を掛ける数列. \booknavref{節：等比型}{本文の説明へ.}
\item[\hypertarget{term:difference}{階差}] 隣り合う項の差.変化を足し戻すと元の数列に戻る. \booknavref{節：階差型}{本文の説明へ.}
\item[\hypertarget{term:ratio}{階比}] 隣り合う項の倍率を追う見方.積による表示と正規化を扱う. \booknavref{節：階比型}{本文の説明へ.}
\item[\hypertarget{term:auxiliary}{補助数列}] 元の数列から,変化を調べやすい別の量を作った数列. \booknavref{節：補助数列の設計}{本文の説明へ.}
\item[\hypertarget{term:normalize}{正規化}] 基準となる量で割るなどして,倍率や大きさをそろえる操作. \booknavref{節：補助数列の設計}{本文の説明へ.}
\item[\hypertarget{term:fixed}{不動点}] 規則を適用しても変わらない値. \booknavref{節：特性方程式型}{本文の説明へ.}
\item[\hypertarget{term:characteristic}{特性方程式}] 一定の値や倍率を使って,元の式を扱いやすくするために解く方程式.本書では不動点を求める式にもこの名称を用いる. \booknavref{節：特性方程式型}{本文の説明へ.}
\item[\hypertarget{term:root}{特性根}] 特性方程式の解. \booknavref{節：二次特性方程式型}{本文の説明へ.}
\item[\hypertarget{term:undetermined}{未定係数法}] 式の形を先に選び,係数を比較して未知の係数を決める方法. \booknavref{節：階差等比複合型}{本文の説明へ.}
\item[\hypertarget{term:iterate}{反復}] 同じ規則を繰り返し適用すること. \booknavref{節：反復・不動点}{本文の説明へ.}
\item[\hypertarget{term:induction}{数学的帰納法}] 出発点と,次の番号へ進めることを示して,すべての番号について証明する方法. \booknavref{節：数学的帰納法}{本文の説明へ.}
\item[\hypertarget{term:linear}{線形}] 入力の和と定数倍を保つ性質.線形漸化式という名称は,未知の項が係数を掛けて足す形で現れる方程式に使う. \booknavref{節：線形・非線形}{本文の説明へ.}
\item[\hypertarget{term:nonlinear}{非線形}] 線形の条件を満たさないもの.項の二乗や逆数を含む例を扱う. \booknavref{節：線形・非線形}{本文の説明へ.}
\item[\hypertarget{term:affine}{アフィン}] 線形な規則に定数を加えた形. \booknavref{節：アフィン・同次・非同次}{本文の説明へ.}
\item[\hypertarget{term:homogeneous}{同次}] 線形な式で,未知の項を集めた後に付加項が残らない形. \booknavref{節：アフィン・同次・非同次}{本文の説明へ.}
\item[\hypertarget{term:inhomogeneous}{非同次}] 線形な式で,未知の項を集めた後に既知の付加項が残る形. \booknavref{節：アフィン・同次・非同次}{本文の説明へ.}
\item[\hypertarget{term:superposition}{重ね合わせの原理}] 同次線形の式を満たす解を,定数倍して足しても解になる性質. \booknavref{節：線形結合}{本文の説明へ.}
\item[\hypertarget{term:combination}{線形結合}] 数列やベクトルに定数を掛けて足したもの. \booknavref{節：線形結合}{本文の説明へ.}
\item[\hypertarget{term:independent}{一次独立}] 定数倍して足したものが零になるのは,係数がすべて零の場合だけ,という関係. \booknavref{節：一次独立・基底・次元}{本文の説明へ.}
\item[\hypertarget{term:basis}{基底}] すべてを定数倍と足し算で表せる,一次独立な組. \booknavref{節：一次独立・基底・次元}{本文の説明へ.}
\item[\hypertarget{term:dimension}{次元}] 基底に必要な対象の個数.解全体では,自由に選べる初期値の個数とつながる. \booknavref{節：一次独立・基底・次元}{本文の説明へ.}
\item[\hypertarget{term:vector-space}{ベクトル空間}] 足し算と定数倍ができ,これらの操作が所定の規則を満たす集合.具体的な解の集合から説明する. \booknavref{節：一次独立・基底・次元}{本文の説明へ.}
\item[\hypertarget{term:constant}{定数係数}] 未知の項に掛かる係数が,添字によらず一定であること. \booknavref{節：定数係数・変数係数}{本文の説明へ.}
\item[\hypertarget{term:variable}{変数係数}] 未知の項に掛かる係数が,添字によって変わること. \booknavref{節：定数係数・変数係数}{本文の説明へ.}
\item[\hypertarget{term:order}{階数}] 次の項を求める際に,どこまで前の項を使うかを表す数.適用する形と範囲は参照先で確認する. \booknavref{節：階数と初期条件}{本文の説明へ.}
\item[\hypertarget{term:elementary}{初等関数}] 多項式・指数・対数・三角関数などを,有限回の演算や合成で組み合わせた関数. \booknavref{節：初等関数・閉じた形}{本文の説明へ.}
\item[\hypertarget{term:closed}{閉じた形}] 番号が増えても式中の演算の個数を増やさずに値を表す言い方.使える関数は文脈による. \booknavref{節：初等関数・閉じた形}{本文の説明へ.}
\item[\hypertarget{term:logarithm}{対数}] ある数が底の何乗かを表す数.真数と底の条件を伴う. \booknavref{節：対数}{本文の説明へ.}
\item[\hypertarget{term:inverse}{逆関数}] 入力と出力の対応を逆にたどる関数. \booknavref{節：対数}{本文の説明へ.}
\item[\hypertarget{term:binomial}{二項定理}] 二つの数の和の整数乗を,項の和へ展開する公式. \booknavref{節：二項定理}{本文の説明へ.}
\item[\hypertarget{term:trigonometric}{三角関数}] 角度を,円上の位置や辺の比と結び付ける関数. \booknavref{節：三角関数}{本文の説明へ.}
\item[\hypertarget{term:addition}{加法定理}] 角度の和を,元の角度の三角関数で表す公式. \booknavref{節：三角関数}{本文の説明へ.}
\item[\hypertarget{term:limit}{極限}] 番号や変数を近づけるときに,値が近づく先を調べる考え方. \booknavref{節：極限}{本文の説明へ.}
\item[\hypertarget{term:one-sided}{片側極限}] 変数を片側から近づけたときの極限. \booknavref{節：片側極限}{本文の説明へ.}
\item[\hypertarget{term:left-limit}{左側極限}] 小さい側から変数を近づけたときの極限. \booknavref{節：片側極限}{本文の説明へ.}
\item[\hypertarget{term:right-limit}{右側極限}] 大きい側から変数を近づけたときの極限. \booknavref{節：片側極限}{本文の説明へ.}
\item[\hypertarget{term:polar}{極形式}] 複素数を,絶対値と角度によって表す形. \booknavref{節：複素平面と極形式}{本文の説明へ.}
\item[\hypertarget{term:complex-plane}{複素平面}] 複素数を平面上の点として表す見方. \booknavref{節：複素平面と極形式}{本文の説明へ.}
\item[\hypertarget{term:conjugate}{共役}] 複素数の虚部の符号を反転させる関係. \booknavref{節：複素平面と極形式}{本文の説明へ.}
\item[\hypertarget{term:argument}{偏角}] 複素平面で,正の実軸から測る角度.角度の範囲の取り方を伴う. \booknavref{節：複素平面と極形式}{本文の説明へ.}
\item[\hypertarget{term:de-moivre}{ド・モアブルの定理}] 複素数の整数乗で,絶対値と角度がどう変わるかを表す関係. \booknavref{節：ド・モアブルの定理}{本文の説明へ.}
\item[\hypertarget{term:euler}{オイラーの公式}] 複素指数関数と正弦・余弦を結ぶ公式. \booknavref{節：オイラーの公式}{本文の説明へ.}
\item[\hypertarget{term:matrix}{行列}] 数を長方形に並べ,複数の量の更新をまとめて表す道具. \booknavref{節：行列の基本}{本文の説明へ.}
\item[\hypertarget{term:column-vector}{列ベクトル}] 複数の数を縦に並べた,一組の記録. \booknavref{節：行列の基本}{本文の説明へ.}
\item[\hypertarget{term:row-vector}{行ベクトル}] 複数の数を横に並べたもの. \booknavref{節：行列の基本}{本文の説明へ.}
\item[\hypertarget{term:identity}{単位行列}] 掛けても元のベクトルや行列を変えない行列. \booknavref{節：行列の基本}{本文の説明へ.}
\item[\hypertarget{term:eigenvector}{固有ベクトル}] 行列を掛けると,元のベクトルの定数倍になる非零のベクトル. \booknavref{節：行列の基本}{本文の説明へ.}
\item[\hypertarget{term:left-eigen}{左固有ベクトル}] 行列を右から掛けると,自身の定数倍になる非零の行ベクトル. \booknavref{節：行列の基本}{本文の説明へ.}
\item[\hypertarget{term:eigenvalue}{固有値}] 固有ベクトルが更新されるときの倍率. \booknavref{節：行列の基本}{本文の説明へ.}
\item[\hypertarget{term:transpose}{転置}] 行と列を入れ替える操作. \booknavref{節：行列の基本}{本文の説明へ.}
\item[\hypertarget{term:invariant}{不変量}] 更新しても値が変わらない量. \booknavref{節：不変量から線形の関係}{本文の説明へ.}
\item[\hypertarget{term:double}{二重数列}] 二つの添字で項を指定する数列. \booknavref{節：二重数列と二方向の階差}{付録の説明へ.}
\item[\hypertarget{term:periodic}{周期性}] 一定の間隔を進めると同じ状態に戻る性質. \booknavref{節：非線形の周期性}{本文の説明へ.}
\item[\hypertarget{term:floor}{ガウス記号}] ある数を超えない最大の整数を表す記号. \booknavref{節：ガウス型}{本文の説明へ.}
\item[\hypertarget{term:floor-function}{床関数}] ある数を超えない最大の整数を返す関数.ガウス記号と結び付く. \booknavref{節：ガウス型}{本文の説明へ.}
\item[\hypertarget{term:newton}{ニュートン法}] 接線の交点から,方程式の解に近い次の値を作る方法. \booknavref{節：ニュートン法型}{本文の説明へ.}
\item[\hypertarget{term:pell}{ペル方程式}] 整数解を求める二次方程式の族.本書では二乗の差が1となる例を追う. \booknavref{節：ペル方程式と近似分数}{本文の説明へ.}
\item[\hypertarget{term:continued}{連分数}] 分母の中に分数が続く表示.近似分数と漸化式を結び付ける. \booknavref{節：ペル方程式と近似分数}{本文の説明へ.}
\item[\hypertarget{term:difference-operator}{差分}] 添字を一つずらした値から,元の値を引く操作. \booknavref{節：差分から始まる離散的な微積分}{本文の説明へ.}
\item[\hypertarget{term:generating}{母関数}] 数列の項を係数に並べ,数列全体を一つの式として扱う道具. \booknavref{節：数列全体を一つの対象にまとめると,何が見えるか}{本文の説明へ.}
\item[\hypertarget{term:convolution}{畳み込み}] 添字の和を合わせた項どうしの積を足す操作.母関数の積に現れる. \booknavref{節：演習の操作を別の表現で見直す}{本文の説明へ.}
\item[\hypertarget{term:chebyshev}{チェビシェフ多項式}] 余弦の整数倍の公式を,多項式の列として表したもの. \booknavref{節：特性方程式を行列で見直す}{本文の説明へ.}
\item[\hypertarget{term:legendre}{ルジャンドル多項式}] 三項漸化式で作られる多項式列.本文では直交性も紹介する. \booknavref{節：特性方程式を行列で見直す}{本文の説明へ.}
\item[\hypertarget{term:interpolation}{ラグランジュ補間}] 有限個の点を通る多項式を作る方法. \booknavref{節：特性方程式を行列で見直す}{本文の説明へ.}
\end{description}

% section: 用語解説

\section{用語の意味}\label{節：用語解説}
  \noindent\textbf{【本編とのつながり】}\par
  式をどの量として見るかを表す用語を確認する.\bookterm{linear}{線形}・\bookterm{nonlinear}{非線形}は\sectionref*{節：漸化式の分類},\bookterm{normalize}{正規化}は\sectionref*{節：階比型},\bookterm{fixed}{不動点}は\sectionref*{節：特性方程式型}へつながる.

  本編や付録では,高校数学では直接習わないような用語も登場する.
  ここでは,\bookterm{recurrence}{漸化式}を読むために役立つ用語を具体例から説明する.
  一度にすべて覚える必要はなく,意味が分からなくなったときに立ち戻ればよい.

  \subsection{補助数列・正規化}\label{節：補助数列・正規化}
    元の\bookterm{sequence}{数列}から,更新を調べやすい別の量を作ったものを\textbf{\bookterm{auxiliary}{補助数列}}という.
    例えば$b_n=a_n-\alpha$や$b_n=\frac1{a_n}$である.逆数や\bookterm{logarithm}{対数}を使う際は,その量が定義できる条件を確かめる.

    本書では,既知の非零の\bookterm{sequence}{数列}$h_n$で割り,倍率や大きさの基準をそろえる操作を\textbf{\bookterm{normalize}{正規化}}と呼ぶ.
    例えば$a_{n+1}=\frac{n+1}{n}a_n$を
    \[
      \frac{a_{n+1}}{n+1}=\frac{a_n}{n}
    \]
    とすると,\,$b_n=\frac{a_n}{n}$は一定になる.
    分母の$n$を$n+1$へ変えた比が,元の倍率$\frac{n+1}{n}$と一致するのが手掛かりである.
    \bookterm{auxiliary}{補助数列}を作る共通の考え方は\sectionref*{節：補助数列の設計}で説明する.
    どの\bookterm{sequence}{数列}で割るかを選ぶ練習は\sectionref*{節：操作を選ぶ練習}で行う.

  \subsection{線形・非線形}\label{節：線形・非線形}
    まず,関数$f$が\textbf{\bookterm{linear}{線形}}であるとは,入力の和と定数倍に対して
    \[
      f(x+y)=f(x)+f(y),\qquad
      f(cx)=cf(x)
    \]
    が成り立つことをいう.ここでは実数の入力と実数の定数$c$を考える.
    例えば$f(x)=2x$は\bookterm{linear}{線形}である.
    一方,$f(x)=x^2$は一般に$(x+y)^2\ne x^2+y^2$なので\bookterm{linear}{線形}ではない.
    このように,\bookterm{linear}{線形}の条件を満たさないものを\textbf{\bookterm{nonlinear}{非線形}}という.

    \bookterm{recurrence}{漸化式}では,未知の\bookterm{sequence}{数列}の項が係数を掛けて足す形だけで現れるものを\textbf{\bookterm{linear}{線形}\bookterm{recurrence}{漸化式}}という.
    例えば
    \begin{align*}
      a_{n+1}&=2a_n+3,\\
      a_{n+2}&=3a_{n+1}-2a_n
    \end{align*}
    は\bookterm{linear}{線形}\bookterm{recurrence}{漸化式}である.
    これに対して
    \[
      a_{n+1}=a_n^2+1,\qquad
      a_{n+1}=\frac{1}{a_n}
    \]
    は\bookterm{nonlinear}{非線形}\bookterm{recurrence}{漸化式}である.後者では$a_n\ne0$が必要である.
    項の2乗,項どうしの積,項の逆数などが現れるかが見分ける目安になる.

    注意したいのは,\bookterm{linear}{線形}かどうかは\textbf{何を未知の量として見ているか}で決まることである.
    \[
      a_{n+1}=n^2a_n+2^n
    \]
    も,$a_n$に関しては\bookterm{linear}{線形}である.
    $n^2$や$2^n$は,\,$n$が決まれば既知の数だからである.
    \bookterm{general}{一般項}のグラフが直線になるという意味ではない.
    実際,$a_{n+1}=2a_n$は\bookterm{linear}{線形}\bookterm{recurrence}{漸化式}だが,その\bookterm{general}{一般項}は$a_n=a_1 2^{n-1}$である.

  \subsection{アフィン・同次・非同次}\label{節：アフィン・同次・非同次}
    $f(x)=px+q$のように,\bookterm{linear}{線形}な関数に定数を加えたものを\textbf{\bookterm{affine}{アフィン}}な関数という.
    $q\ne0$なら$f(0)\ne0$であり,関数としては\bookterm{linear}{線形}ではない.
    高校数学の一次関数は,$p\ne0$の\bookterm{affine}{アフィン}な関数に当たる.
    それでも$a_{n+1}=pa_n+q$を\bookterm{linear}{線形}\bookterm{recurrence}{漸化式}と呼ぶのは,未知の項について\bookterm{linear}{線形}な式に,既知の項を加えた方程式だからである.

    \bookterm{linear}{線形}\bookterm{recurrence}{漸化式}で,未知の項をすべて左辺に集めて
    \[
      a_{n+1}-pa_n=q
    \]
    と書いたとき,右辺が$0$なら\textbf{\bookterm{homogeneous}{同次}},そうでなければ\textbf{\bookterm{inhomogeneous}{非同次}}という.
    「\bookterm{homogeneous}{同次}」は「斉次」とも書く.
    右辺は定数に限らず,$n$についての既知の式でもよい.
    例えば$a_{n+1}-2a_n=n$は\bookterm{inhomogeneous}{非同次}である.

    \bookterm{homogeneous}{同次}\bookterm{linear}{線形}\bookterm{recurrence}{漸化式}では,2つの解の\bookterm{sequence}{数列}$u_n,v_n$から,定数$C,D$を用いて作る$Cu_n+Dv_n$も解になる.
    この性質を\textbf{\bookterm{superposition}{重ね合わせの原理}}という.
    例えば$u_{n+1}=2u_n$,\,$v_{n+1}=2v_n$なら
    \[
      Cu_{n+1}+Dv_{n+1}=2(Cu_n+Dv_n)
    \]
    となる.ただし,同じ\bookterm{initial}{初期条件}を満たすとは限らない.
    また,\bookterm{inhomogeneous}{非同次}の場合は一般にこの性質をそのまま使えない.
    この性質から解全体を見渡す\bookterm{combination}{線形結合}・\bookterm{independent}{一次独立}・\bookterm{dimension}{次元}の考え方へ進む道筋は,
    \sectionref*{節：解を組み合わせる見方}で具体例から説明する.

    \bookterm{inhomogeneous}{非同次}の式を\bookterm{homogeneous}{同次}の式に直す例も見ておこう.
    $a_{n+1}=2a_n+3$に対して$b_n=a_n+3$と置くと,
    \[
      b_{n+1}=a_{n+1}+3=2(a_n+3)=2b_n
    \]
    となる.こうした置き換えが,本編で扱う解法の基本になる.

  \subsection{定数係数・変数係数}\label{節：定数係数・変数係数}
    \bookterm{linear}{線形}\bookterm{recurrence}{漸化式}で未知の項に掛かる数を\textbf{係数}という.
    係数が$n$によらず一定なら\textbf{\bookterm{constant}{定数係数}},\,$n$によって変わるなら\textbf{\bookterm{variable}{変数係数}}という.
    例えば
    \[
      a_{n+2}=3a_{n+1}-2a_n
    \]
    は\bookterm{constant}{定数係数}であり,
    \[
      a_{n+1}=(n+1)a_n
    \]
    は\bookterm{variable}{変数係数}である.
    $a_{n+1}=2a_n+n^2$では,未知の項の係数は一定なので\bookterm{constant}{定数係数}である.
    右辺の$n^2$は,未知の項に掛かる係数ではなく\bookterm{inhomogeneous}{非同次}の項である.
    \bookterm{constant}{定数係数}かどうかと,\bookterm{homogeneous}{同次}かどうかは別々に判断する.

  \subsection{階数と初期条件}\label{節：階数と初期条件}
    \bookterm{recurrence}{漸化式}の\textbf{\bookterm{order}{階数}}は,関係する項の添字の最大値と最小値の差をいう.
    例えば$a_{n+1}=2a_n$は\textbf{一階},
    \[
      a_{n+2}=a_{n+1}+a_n
    \]
    は\textbf{二階}の\bookterm{recurrence}{漸化式}である.
    二階は「隣接3項間」に対応するが,$a_{n+2}=a_n$も二階である.
    また,$a_{n+1}=a_n^2$は\bookterm{nonlinear}{非線形}だが一階であり,\bookterm{order}{階数}と項のべきの大きさは異なる.

    \bookterm{sequence}{数列}を具体的に決めるために最初に与える値を\textbf{\bookterm{initial}{初期条件}}という.
    次の項が直前の$k$項から一意に計算できる$k$階の\bookterm{recurrence}{漸化式}では,通常$a_1,\ldots,a_k$を指定する.
    上の二階の例なら$a_1,a_2$から$a_3,a_4,\ldots$を順に求められる.
    分母が$0$になるなど,更新の途中で式が定義されなくなる場合には別途注意が必要である.

  \subsection{反復・不動点}\label{節：反復・不動点}
    $a_{n+1}=f(a_n)$では,同じ関数$f$を繰り返し適用して\bookterm{sequence}{数列}を作る.
    この操作を\textbf{\bookterm{iterate}{反復}}という.
    \[
      a_2=f(a_1),\qquad a_3=f(f(a_1))
    \]
    となり,$f(f(x))$は$f(x)$の2乗ではない.

    $f(\alpha)=\alpha$を満たす値$\alpha$を\textbf{\bookterm{fixed}{不動点}}という.
    その値から出発すれば,更新しても値が変わらない.
    例えば$f(x)=2x+3$の\bookterm{fixed}{不動点}は
    \[
      \alpha=2\alpha+3\iff\alpha=-3
    \]
    である.$a_1=-3$ならすべての項が$-3$になる.
    一般の初項でも,\bookterm{fixed}{不動点}からのずれ$b_n=a_n-\alpha$を考えると
    \[
      b_{n+1}=2b_n
    \]
    と簡単になる.先ほどの$b_n=a_n+3$という置き換えは,\bookterm{fixed}{不動点}からのずれを見ていたのである.
    \bookterm{fixed}{不動点}があっても,そこへ収束するとは限らない.
    この例では$a_1\ne-3$なら$|b_n|$は一歩ごとに2倍になる.

  \subsection{初等関数・閉じた形}\label{節：初等関数・閉じた形}
    多項式,有理関数,指数関数,\bookterm{logarithm}{対数}関数,\bookterm{trigonometric}{三角関数},逆\bookterm{trigonometric}{三角関数}などを,有限回の四則演算や合成で組み合わせて得られる関数を\textbf{\bookterm{elementary}{初等関数}}という.
    合成とは,例えば$\log(1+x^2)$のように,ある関数の値を別の関数に入力することである.
    各関数や演算が定義される範囲で考える.

    \bookterm{general}{一般項}を\textbf{\bookterm{closed}{閉じた形}}で表すという言い方もある.
    何を使ってよいかは文脈によるが,本書では\bookterm{elementary}{初等関数}による表現を一つの目安にしている.
    例えば$a_1=0$,\,$a_{n+1}=a_n+n$から
    \[
      a_n=\sum_{k=1}^{n-1}k=\frac{n(n-1)}{2}
    \]
    が得られる.最後の式なら,$n$が大きくなっても足し算の回数を増やさずに値を求められる.
    一方,$\sum_{k=1}^{n-1}k$も\bookterm{general}{一般項}の表現ではあるが,和に含まれる項数は$n$とともに増える.
    \bookterm{closed}{閉じた形}が見つからなくても,項を順に計算したり,\bookterm{sequence}{数列}の性質を調べたりすることはできる.

% section: シグマの計算

\section{シグマの計算}\label{節：シグマの計算}
  \noindent\textbf{【本編とのつながり】}\par
  和の記号と添字の役割を確認する.\sectionref*{節：階差型}で変化を足し戻し,\sectionref*{節：総和型}で総和の式をずらして引くときに使う.

  $1^2+2^2+3^2+4^2$のような和は,
  \[
    \sum_{k=1}^{4}k^2=1^2+2^2+3^2+4^2=30
  \]
  と短く書ける.$k$を$1$から$4$まで一つずつ動かし,そのたびに$k^2$を足すという意味である.
  この$\sum$を\bookterm{sigma}{シグマ}という.
  実数の\bookterm{sequence}{数列}$\{a_k\}$\footnote{複素数でも構わない.}に対しても,\,$a_1$から$a_n$までの総和を
  \[
    \sum_{k=1}^{n} a_k=a_1+a_2+\cdots+ a_n
  \]
  と表す.
  添字$k$は,和の中で順に動く番号である.同じ範囲なら$\sum_{k=1}^{n}a_k=\sum_{j=1}^{n}a_j$と書き換えても値は変わらない.
  一方,上端の$n$は足す項数を決める.例えば$S_n=\sum_{k=1}^{n}a_k$なら
  \[
    S_{n+1}=S_n+a_{n+1}
  \]
  である.
  足す項がない空の和は$0$とする.例えば$\sum_{k=1}^{n-1}a_k$は,$n=1$のとき$0$である.
  この時以下の事柄が一般に成り立つ.
  ここで$c$は,和を取る添字$k$に依存しない定数である.
  \begin{align*}
    &\sum_{k=1}^{n} c=nc\\
    &\sum_{k=1}^{n} ca_k=c\sum_{k=1}^{n} a_k\\
    &\sum_{k=1}^{n} (a_k+b_k)=\sum_{k=1}^{n} a_k+\sum_{k=1}^{n} b_k
  \end{align*}
  「定数」は,ここでは$k$を動かしても変わらないという意味である.
  例えば$\sum_{k=1}^{n}nk=n\sum_{k=1}^{n}k$では,$n$を和の外へ出せる.$k$そのものは外へ出せない.
  また,一般に以下が成り立つ.
  \begin{align*}
    &\sum_{k=1}^{n} k=\frac{1}{2}n(n+1)\\
    &\sum_{k=1}^{n} k^2=\frac{1}{6}n(n+1)(2n+1)\\
    &\sum_{k=1}^{n} k^3=\left\{\frac{1}{2}n(n+1)\right\}^2\\
  \end{align*}

  \bookterm{recurrence}{漸化式}の添字をずらす場合は,式全体の$n$を置き換える.
  例えば$a_{n+1}=2a_n+n$から$a_{n+2}=2a_{n+1}+n+1$を得る.

% section: 式の変形

\section{式の変形}\label{節：式の変形}
  \noindent\textbf{【本編とのつながり】}\par
  同じ値を保ちながら既知の形を作る操作を確認する.定数の調整は\sectionref*{節：特性方程式型},\bookterm{partial}{部分分数分解}は\sectionref*{節：階差型}で和を計算する際に使う.

  \subsection{同じ数を足して引く}\label{節：+a-aをする}
    $x^2+6x+5$を,一つの2乗を使って表したい.
    $(x+3)^2=x^2+6x+9$なので,$x^2+6x$の部分に$9$を足せば平方になる.値を変えないため,同じ$9$を引く.
    \[
      \begin{aligned}
        x^2+6x+5&=x^2+6x+9-9+5\\
        &=(x+3)^2-4.
      \end{aligned}
    \]
    値を変えずに,$x$を含む部分を一つの2乗にまとめられた.
    この操作を\bookterm{square}{平方完成}という.
    一般の二次式でも,$a\ne0$として同じように進める.
    \begin{align*}
      &ax^2+bx+c\\
      &\qquad=a\left(x^2+\frac{b}{a}x\right) +c\\
      &\qquad=a\left(x^2+\frac{b}{a}x+\frac{b^2}{4a^2}-\frac{b^2}{4a^2}\right)+c \booktag{square-adjust}{☆}\\
      &\qquad=a\left(x+\frac{b}{2a}\right)^2-\frac{b^2}{4a}+c
    \end{align*}
    と,$x$に関する項を1つにまとめられる.
    ここで注目してほしいのは\bookeqbackref{square-adjust}{☆}部分の式変形である.
    これが意味するのは,無理やりに2乗の形,すなわち平方を作り上げよう,ということだ.
    これにより$x$に関する項は定数とともに括弧の中に閉じ込められ,よりシンプルに式が記述できる.
    二次関数の最大・最小問題や軸の決定などにおいても\bookterm{square}{平方完成}は有用である.

    ここから学ぶべきことは,定数項のやり取りによって式を無理やりわかりやすい形,知っている形,扱いやすい形に変形できるということだ.

  \subsection{部分分数分解}\label{節：部分分数分解}
    今度は,一つの分数を簡単な分数の和に分けてみよう.
    例えば,$\dfrac1{x+1}-\dfrac1{x+2}$を通分すると,
    \[
      \begin{aligned}
        \frac1{x+1}-\frac1{x+2}
          &=\frac{(x+2)-(x+1)}{(x+1)(x+2)}\\
          &=\frac1{(x+1)(x+2)}
      \end{aligned}
    \]
    となる.ここでは$x\ne-1,-2$とする.
    この等式を逆向きに使えば,分母が二次式の分数を,分母が一次式の二つの分数へ分けられる.
    $x$を一つずつ増やして和を取ると,隣の項と打ち消し合う形でもある.
    このような操作を\bookterm{partial}{部分分数分解}という.

    一般に
    \[
      \frac{1}{(x+a)(x+b)}
    \]
    でも,分母を一次式に分けることを考える.
    以下では$x\ne-a,-b$とし,$a\ne b$の場合を扱う.
    \[
      \frac{1}{(x+a)(x+b)}=\frac{\alpha}{x+a}+\frac{\beta}{x+b}
    \]
    分母が零にならない\textbf{すべての$x$で一致する}ように,$\alpha,\beta$を決める.
    右辺を通分すると,
    \begin{align*}
    \frac{1}{(x+a)(x+b)}
      &= \frac{\alpha}{x+a}+\frac{\beta}{x+b}\\
      &= \frac{(\alpha+\beta)x+\alpha b+\beta a}
      {(x+a)(x+b)}
    \end{align*}
    分子が常に$1$になるには,$x$の係数が$0$,定数項が$1$であればよい.
    \[
    \therefore\quad
    \begin{cases}
    \alpha+\beta=0,\\
    \alpha b+\beta a=1
    \end{cases}
    \]
    これを解くと,\,$\beta=-\alpha$であるから,
    \[
      \alpha b-\alpha a=1
      \iff \alpha(b-a)=1
    \]
    したがって,\,$a\ne b$であれば,
    \[
      \alpha=\frac{1}{b-a},
      \qquad
      \beta=-\frac{1}{b-a}
    \]
    が得られる.従って
    \[
      \frac1{(x+a)(x+b)}
      =\frac1{b-a}\left(\frac1{x+a}-\frac1{x+b}\right).
    \]
    $a=b$のときは分母に同じ因子が二つ現れるため,この形の分解はできない.
    3次の場合も,\,$p,q,r\ne0$で,分母の一次因子$px+a,qx+b,rx+c$の根
    \[
      -\frac{a}{p},\qquad -\frac{b}{q},\qquad -\frac{c}{r}
    \]
    が互いに異なるときは,
\BookMathLayout{\[
      \frac{1}{(px+a)(qx+b)(rx+c)}=\frac{\alpha}{px+a}+\frac{\beta}{qx+b}+\frac{\gamma}{rx+c}
    \]}{
\[
\begin{aligned}
\frac{1}{(px+a)(qx+b)(rx+c)}
&=\frac{\alpha}{px+a}+\frac{\beta}{qx+b}\\
&\quad+\frac{\gamma}{rx+c}
\end{aligned}
\]
}
    と置き,係数$\alpha,\beta,\gamma$を定めることで分解できる.
    言い換えれば,これらの一次因子が互いに定数倍でないことが条件である.
    重複因子があるときは,その重複度に応じて例えば$\dfrac{A}{x+a}+\dfrac{B}{(x+a)^2}$のような項も必要になる.

% section: 等差数列とその総和の導出

\section{等差数列の和}\label{節：等差数列とその総和の導出}
  \sectionref*{節：階差型}で一次式を足すときに使う.
  初項$a$,公差$d$,末項$l=a+(n-1)d$の\bookterm{arithmetic}{等差数列}の和を$S_n$とする.

  \subsection{和の公式}\label{節：等差数列とその総和の導出：総和の公式}
    \[
      S_n=\frac n2(a+l)=\frac n2\{2a+(n-1)d\}.
    \]
    項数と,初項・末項の平均との積である.

  \subsection{公式の導出}\label{節：等差数列とその総和の導出：公式の導出}
    $d>0$の場合,増分$d,2d,\ldots,(n-1)d$を逆順の列と重ねると,各列の高さが$nd$にそろう.
    \[
      2\sum_{k=1}^{n-1}kd=n(n-1)d
    \]
    に初項$a$を$n$回加えれば,$S_n=na+n(n-1)d/2$を得る.
    \BookFigArithmeticSum

    同じ工夫を元の数列に使うと,$d$の符号によらず導ける.
    \[
      \begin{array}{rcl}
        S_n&=&a+(a+d)+(a+2d)\\
           &&\qquad+\cdots+\{a+(n-1)d\},\\
        S_n&=&\{a+(n-1)d\}+\{a+(n-2)d\}\\
           &&\qquad+\cdots+(a+d)+a.
      \end{array}
    \]
    各列の和は$2a+(n-1)d$であり,$n$列あるので$2S_n=n\{2a+(n-1)d\}$となる.

% section: 等比数列の総和とその導出

\section{等比数列の和}\label{節：等比数列の総和とその導出}
  \noindent\textbf{【本編とのつながり】}\par
  \bookterm{geometric}{等比数列}の和を,倍率をそろえて引く操作から導く.\sectionref*{節：階差型}と\sectionref*{節：特性方程式型}の\bookterm{difference}{階差}による別解で使う.

  初項が$a$,公比が$r$である\bookterm{geometric}{等比数列}
  \[
    a,\,ar,\,ar^2,\,\ldots,\,ar^{n-1}
  \]
  の初項から第$n$項までの総和を考える.
  その総和を$S_n$とすると,
  \[
    S_n
    =\sum_{k=1}^{n} ar^{k-1}
    =a+ar+ar^2+\cdots+ar^{n-1}
  \]
  と表せる.

  \subsection{和の公式}\label{節：等比数列の総和とその導出：総和の公式}
    公比$r$が$1$でないとき,
    \[
      S_n=\frac{a(1-r^n)}{1-r}
    \]
    である.
    分母と分子の符号を同時に変えれば,
    \[
      S_n=\frac{a(r^n-1)}{r-1}
    \]
    とも書ける.
    公比が$1$のときは,すべての項が$a$であるから,
    \[
      S_n=na
    \]
    となる.

  \subsection{公式の導出}\label{節：等比数列の総和とその導出：公式の導出}
    $r\ne1$として,総和の式の両辺に$r$を掛ける.
    \begin{align*}
      S_n
        &=a+ar+ar^2+\cdots+ar^{n-1},\\
      rS_n
        &=ar+ar^2+ar^3+\cdots+ar^n.
    \end{align*}
    2つの式を項ごとに引くと,右辺では共通する項がすべて消える.
\BookMathLayout{\[
    \begin{array}{rrrl}
        &S_{n}={}&a+&\hspace{-0.75em}ar+ar^2+\cdots+ar^{n-1}\\
      -\,) & rS_n={}&{}&\hspace{-0.75em}ar+ar^2+ar^3+\cdots+ar^{n}\\ \hline
        & (1-r)S_n={}&{}&\hspace{-0.75em}a-ar^n
    \end{array}
    \]}{
\[
\begin{array}{rl}
S_n={}&a+ar+ar^2+\cdots\\
     &{}+ar^{n-1}\\
-\,)\ rS_n={}&ar+ar^2+ar^3+\cdots\\
     &{}+ar^n\\ \hline
(1-r)S_n={}&a-ar^n
\end{array}
\]
}
    よって,
    \[
      (1-r)S_n=a(1-r^n)
    \]
    である.
    $r\ne1$より両辺を$1-r$で割ると,
    \[
      S_n=\frac{a(1-r^n)}{1-r}
    \]
    を得る.

    この導出では,総和を公比$r$倍してから元の総和との差を取ることで,途中の項を消している.
    \bookterm{sequence}{数列}の和を求めるときには,このように総和を定数倍して引き,共通する項を消す考え方がよく用いられる.

% section: 数学的帰納法

  \section{数学的帰納法}\label{節：数学的帰納法}
  \noindent\textbf{【本編とのつながり】}\par
  数項からの予想をすべての項の証明へ進める方法を学ぶ.\sectionref*{節：漸化式を解く前に}の検算や,\sectionref*{節：虚数解型隣接3項間漸化式}の\bookterm{general}{一般項}の確認で使う.

    ここまでは計算に用いる公式や式変形を扱った. ここでは,自然数に関する主張を証明する方法を説明する.
    自然数$n$についての主張$P(n)$が,すべての$n\ge1$で成り立つことを示したいとする.
    このとき,以下の2つを示せばよい.
    \begin{enumerate}[label=(\roman*),parsep=3pt]
      \item $P(1)$が成り立つ.
      \item $n\ge1$を任意にとり,\,$P(n)$が成り立つと仮定すると,\,$P(n+1)$も成り立つ.
    \end{enumerate}
    この2つが示されれば,\,(i)より$P(1)$が成り立ち,\,(ii)より$P(1)\Rightarrow P(2)$,\,$P(2)\Rightarrow P(3)$,\,$\ldots$と次々に成り立つことが分かるので,すべての$n\ge1$について$P(n)$が成り立つ.
    ドミノ倒しに例えられることが多い.\,(i)は最初のドミノを倒すこと,\,(ii)は「どのドミノについても,それが倒れれば次のドミノも倒れる」ように並べることに対応する.

    \sectionref*{節：追う量を変える}の「発見と証明」で確かめた$a_n=2n-1$は,この二段階の具体例である.
    初項の一致が(i),$a_k=2k-1$を仮定して更新式から$a_{k+1}=2(k+1)-1$を得た計算が(ii)に当たる.
    すべての$n$で正しいと先に仮定するのではなく,\textbf{ある番号で正しければ次でも正しい}ことを示している.

    \sectionref*{節：極形式と複素平面}では,この方法を用いて\bookterm{de-moivre}{ド・モアブルの定理}を証明する.
    また,\,\chapterref{章：漸化式の解説}では,\bookterm{recurrence}{漸化式}から予想した\bookterm{general}{一般項}の形を確かめる場面や,\bookterm{trigonometric}{三角関数}の\bookterm{addition}{加法定理}から\bookterm{recurrence}{漸化式}を導く場面(\sectionref{節：漸化式を解く前に}や\sectionref{節：虚数解型隣接3項間漸化式}など)で\bookterm{induction}{数学的帰納法}を用いる.

% section: 指数対数

  \section{指数と対数}\label{節：指数対数}
  \noindent\textbf{【本編とのつながり】}\par
  指数法則と,積・商・累乗を和・差・定数倍に変える\bookterm{logarithm}{対数}法則を確認する.\sectionref*{節：対数型}では,真数の正値と底の条件も使う.

    \subsection{指数}\label{節：指数}
      $a$が$p$回掛けられている状況を考える.
      ここで
      \[
        a\times a\times\cdots\times a
      \]
      のように記すのは煩雑なので,
      \[
        a\times a\times\cdots\times a=a^p
      \]
      と表す.
      このときの$p$を指数といい,
      \[
        f(x)=a^x
      \]
      で表されるような関数を指数関数という.
      ここでは$p$を正の整数として導入した.底$a>0$の場合は,\,$a^0=1,\ a^{-p}=\dfrac1{a^p}$と定め,さらに有理数・実数の指数へと拡張できる.
      以降,実数の指数を用いるときは底$a>0$とする.
    \subsection{対数}\label{節：対数}
      $a^p=q$のとき,\,$p$は
      \[
        \log_a{q}=p
      \]
      のように表される.
      この時$\log$をログと読み,\,$p$を\bookterm{logarithm}{対数}という.
      また,\,$a$を底,\,$q$を真数と言う.
      一般に底$a$は$a>0,\,a\ne1$を満たす必要がある.
      また,実数の\bookterm{logarithm}{対数}では真数$q>0$とする.
      また,
      \[
        e=\lim_{x \to \infty} \left(1+\frac{1}{x}\right)^x
      \]
      で定義される数$e$を自然\bookterm{logarithm}{対数}の底またはネイピア数という.
      ただし,以下の公式では底$a,c$は正で$1$ではなく,真数$p,M,N,b$は正とする.
      \bookterm{logarithm}{対数}に関して,一般に以下が成り立つ.
      \begin{align*}
        &\log_a{p^q}=q\log_a{p}\\
        &\log_a{MN}=\log_a{M}+\log_a{N}\\
        &\log_a{\frac{M}{N}}=\log_a{M}-\log_a{N}\\
        &\log_a{b}=\frac{\log_c{b}}{\log_c{a}}
      \end{align*}

      また,
      \[
        f(x)=\log_a x
      \]
      というような関数を\bookterm{logarithm}{対数}関数といい,指数関数の\bookterm{inverse}{逆関数}である.
      \bookterm{inverse}{逆関数}を作れるのは,出力から元の入力を一つに決められるときである.
      関数$f$の入力の範囲を$D$,実際に取る値の範囲を$E$とする.
      $D$を\textbf{定義域},$E$を\textbf{値域}という.
      各$y\in E$について
      \[
        y=f(x)
      \]
      を満たす$x\in D$がただ一つなら,その$x$を返す関数
      \[
        x=g(y)
      \]
      を作れる.これが$f$の\bookterm{inverse}{逆関数}であり,定義域は$E$,値域は$D$になる.
      元のグラフ上の点$(x,y)$は逆関数のグラフ上の点$(y,x)$に対応するので,二つのグラフは直線$y=x$に関して対称である.
      関数$g$はよく,\,$f^{-1}$と書かれる.
      例えば$x\mapsto a^x$は実数全体から正の実数全体への対応であり,正の値$q$から指数$x$を一つに戻せる.
      その逆が$q\mapsto\log_a q$である.
      一方,$f(x)=x^2$を実数全体で考えると,$4$から$2$と$-2$のどちらにも戻れてしまう.
      定義域を$x\ge0$に限れば一つに決まり,逆関数は$g(y)=\sqrt y$となる.
      このように,逆関数では式だけでなく入力の範囲も決めておく必要がある.\footnote{$f^{-1}(y)$は逆関数の値を表す記号であり,$1/f(y)$ではない.}

% section: 二項定理

  \section{二項定理}\label{節：二項定理}
  \noindent\textbf{【本編とのつながり】}\par
  和の累乗を展開する公式と組合せの意味を確認する.\sectionref*{節：極形式と複素平面}の\bookterm{euler}{オイラーの公式}の導出や,\sectionref*{節：発展難易度}の数え上げで使う.

    次に,和の累乗を展開するための公式を確認する.
    非負の整数$\,n$と実数$\,x,y$について,
    \[
      (x+y)^n=\sum_{k=0}^{n}\binom{n}{k}x^{n-k}y^{k}=\sum_{k=0}^{n}{}_n\mathrm{C}_kx^{n-k}y^{k}
    \]
    が成り立つ.これを\bookterm{binomial}{二項定理}という.\footnote{この多項式の式では,0乗は因子を含まない積を表すので,底が$0$でも$1$と読む.}
    ここで
\BookMathLayout{\[
      {}_n\mathrm{C}_k=\binom{n}{k}
      =\frac{n!}{k!(n-k)!}
      =\frac{n(n-1)\cdots(n-k+1)}{k!}
    \]}{
\[
\begin{aligned}
{}_n\mathrm C_k&=\binom nk=\frac{n!}{k!(n-k)!}\\
&=\frac{n(n-1)\cdots(n-k+1)}{k!}
\end{aligned}
\]
}
    が成り立つ.\,$\binom{n}{k}$および${}_n\mathrm{C}_k$は$\,n$個から$\,k$個を選ぶ組合せの数を表す.
    この階乗表示を使うのは,整数$k$が$0\le k\le n$を満たすときである.\,$0!=1$とする.
    $k=0$では,一つも選ばない選び方が1通りあるので$\binom n0=1$となる.
    また,本書では整数$k$が$k<0$または$k>n$を満たすとき,
    \[
      \binom nk=0
    \]
    と定める.その個数を選ぶ方法がないことに対応している.
    この場合に,階乗表示へ負の整数を代入してはいけない.
    例えば$\binom12=0$であり,
    \[
      \binom n2=\frac{n(n-1)}2\qquad(n\ge0)
    \]
    は$n=0,1$でも使える.
    \booknavref{節：総和型}{総和型}で現れた$n(n-1)/2$も,この書き方を使っている.
    \booknavref{節：二項係数と差分}{二項係数と差分}では,差を取ると添字が一つ下がることを確かめる.

    特に,\bookterm{binomial}{二項定理}ではこの組合せの数を二項係数と呼ぶ.
    その理由を,因子を選ぶ方法から考える.\,$(x+y)^n$は,\,$\,n$個の因子
    \[
      (x+y)(x+y)\cdots(x+y)
    \]
    の積である.展開するとき,各因子から$x$か$y$のどちらか一方を選ぶ.
    このとき,\,$\,x^{n-k}y^k$を作るには,\,$\,n$個の因子のうち$\,k$個から$y$を選び,
    残りの$\,n-k$個から$x$を選べばよい.
    $y$を選ぶ因子の位置の選び方は,\,$\,n$個から$\,k$個を選ぶ組合せの数
    $\binom{n}{k}$通りである.したがって,\,$\,x^{n-k}y^k$の係数は$\binom{n}{k}$となる.
    例えば$\,(x+y)^3$では,\,$\,x+y$という3個の因子のうち,
    どの2個から$x$を選ぶかによって$x^2y$が生じる.
    その選び方は$\binom{3}{2}=3$通りなので,\,$\,x^2y$の係数は$3$である.
    同様に,\,$\,xy^2$の係数も$\binom{3}{1}=3$となり,
    \[
      (x+y)^3=x^3+3x^2y+3xy^2+y^3
    \]
    を得る.
    また,これはパスカルの三角形の3段目の係数と一致する.
    複素数に対しても,同じ\bookterm{binomial}{二項定理}が成り立つ.

\section{三角関数}\label{節：三角関数}
  \noindent\textbf{【本編とのつながり】}\par
  角度と座標の対応から,\bookterm{addition}{加法定理}と倍角公式を確認する.\sectionref*{節：三角関数型}で公式の形を手掛かりにし,\sectionref*{節：虚数解型隣接3項間漸化式}で実数の\bookterm{general}{一般項}を表すために使う.

  原点を中心とする半径$1$の円を考える.
  正の横軸から反時計回りに角度$x$だけ進んだ点の横座標が$\cos x$,縦座標が$\sin x$である.
  時計回りの角度は負に数える.
  弧度法では,この円で進んだ弧の長さを角度とし,一周は$2\pi$となる.
  したがって,
  \[
    \cos^2x+\sin^2x=1
  \]
  である.また,\,$\cos x\ne0$のとき,
  \[
    \tan x=\frac{\sin x}{\cos x}
  \]
  と定める.これは,原点とその点を結ぶ直線の傾きに当たる.
  角度を一周増やしても座標は変わらないので,\,$\sin x$と$\cos x$は$2\pi$の周期をもつ.
  \booknavref{節：複素平面と極形式}{複素数を絶対値と角度で表す方法}も,この座標の対応を使う.

  \bookstep{角度を足す公式}
  \bookterm{trigonometric}{三角関数}には,角度の和を表す\bookterm{addition}{加法定理}
  \begin{align*}
    \sin(\alpha+\beta)
      &=\sin\alpha\cos\beta+\cos\alpha\sin\beta\\
    \cos(\alpha+\beta)
      &=\cos\alpha\cos\beta-\sin\alpha\sin\beta\\
    \tan(\alpha+\beta)
      &=\frac{\tan\alpha+\tan\beta}
              {1-\tan\alpha\tan\beta}
  \end{align*}
  が成り立つ.\,$\sin$と$\cos$の式は,すべての実数$\alpha,\beta$で使える.
  以下の$\tan$を含む式は,両辺が定義される範囲で使う.
  例えば上の式では,$\cos\alpha$,\,$\cos\beta$,\,$\cos(\alpha+\beta)$がいずれも零でないことを確かめる.
  \bookstep{同じ角度を二度,三度足す}
  ここで $\alpha=\beta=x$ とすれば,倍角公式
  \begin{align*}
    \sin 2x&=2\sin x\cos x\\
    \cos 2x&=2\cos^2x-1=1-2\sin^2x\\
    \tan 2x&=\frac{2\tan x}{1-\tan^2x}
  \end{align*}
  を得る.また,倍角公式から
  \begin{align*}
    \sin^2x&=\frac{1-\cos 2x}{2}\\
    \cos^2x&=\frac{1+\cos 2x}{2}\\
    \tan^2x&=\frac{1-\cos 2x}{1+\cos 2x}
  \end{align*}
  を導くことができる.また,\bookterm{addition}{加法定理}から,\,$3x=2x+x$とすることで,
  \begin{align*}
    \sin 3x&=3\sin x-4\sin^3x\\
    \cos 3x&=4\cos^3x-3\cos x
  \end{align*}
  が求められる.

% section: 極限

\section{極限}\label{節：極限}
  \noindent\textbf{【本編とのつながり】}\par
  \bookterm{sequence}{数列}や関数が近づく値を表す記法と計算を確認する.\sectionref*{節：ニュートン法型},発展演習の\bookterm{limit}{極限}問題,\chapterref{章：付録}で使う.基本・応用パターンの\bookterm{general}{一般項}を求める段階では後回しにできる.

  \subsection{極限の考え方}\label{節：極限：極限}
    ある関数$f(x)$について\,$x$を$a$に限りなく近づける操作を数学的には
    \[
      \lim_{x\to a} f(x)
    \]
    と表す.
    これを用いることで,ある関数$f(x)$の$x=a$の近傍での値がどうなっているかを知ることができる.
    例えば,\,$f(x)=x^2-4$のとき
    \[
      \lim_{x\to 2} (x^2-4)
    \]
    は$0$となる.
    一見シンプルな代入計算に見えるかもしれないが,これが今後重要な役割を果たす.

  \subsection{数列の極限}\label{節：数列の極限}
    \bookterm{sequence}{数列}$\{a_n\}$の番号$n$を大きくしていくとき,\,$a_n$がある値$L$に限りなく近づくなら,\,$a_n$は$L$に収束するといい,
    \[
      \lim_{n\to\infty}a_n=L
    \]
    と書く.

    例えば,\,$n>0$のとき
    \[
      a_n=\frac{1}{n}
    \]
    とすると,\,$n$を大きくするほど$a_n$は$0$に近づくので,
    \[
      \lim_{n\to\infty}\frac{1}{n}=0
    \]
    である.

  \subsection{関数の極限}\label{節：関数の極限}
    $x$を$a$に近づけたとき,\,$f(x)$が$L$に近づくなら,
    \[
      \lim_{x\to a}f(x)=L
    \]
    と書く.
    このとき,\,$x=a$での関数の値そのものではなく,\,$a$の近くでの値の変化に注目する.

    \bookterm{limit}{極限}の計算では,既に求めた\bookterm{limit}{極限}を組み合わせたり,不等式で上下から挟んだりする.
    例えば,
    \[
      a_n\le b_n\le c_n
    \]
    が成り立っており,
    \[
      \lim_{n\to\infty}a_n=\lim_{n\to\infty}c_n=L
    \]
    であるとき,\,$b_n$も$L$に収束する.
    これをはさみうちの原理という.

  \subsection{片側極限}\label{節：片側極限}
    $x$を$a$に近づけるとき,近づく側を限定した\bookterm{limit}{極限}を\textbf{\bookterm{one-sided}{片側極限}}という.
    $x<a$の範囲で$a$に近づける場合を\textbf{\bookterm{left-limit}{左側極限}},\,$x>a$の範囲で$a$に近づける場合を\textbf{\bookterm{right-limit}{右側極限}}という.
    数直線で考えると,それぞれ$a$の左側・右側から近づけることに対応する.

    \bookterm{left-limit}{左側極限}には,次のような複数の記法がある.
    \[
      \boxed{
        \begin{aligned}
          \lim_{x\to a-}f(x)
          &=\lim_{x\to a^-}f(x)\\
          &=\lim_{\substack{x\to a\\x<a}}f(x)\\
          &=\lim_{x\uparrow a}f(x)
        \end{aligned}
      }
    \]
    これらは同じ\bookterm{left-limit}{左側極限}を表す記法である.
    $a^-$の$-$は,\,$x<a$の側から近づけることを示す.
    また,\,$x\uparrow a$も$a$より小さい側から$a$に近づけることを表す.

    同様に,\bookterm{right-limit}{右側極限}には次の記法がある.
    \[
      \boxed{
        \begin{aligned}
          \lim_{x\to a+}f(x)
          &=\lim_{x\to a^+}f(x)\\
          &=\lim_{\substack{x\to a\\x>a}}f(x)\\
          &=\lim_{x\downarrow a}f(x)
        \end{aligned}
      }
    \]
    $a^+$の$+$は,\,$x>a$の側から近づけることを示す.
    また,\,$x\downarrow a$も$a$より大きい側から$a$に近づけることを表す.
    特に$a=0$のときは,\,$x\to0^+$を$x\to+0$,\,$x\to0^-$を$x\to-0$と書くこともある.
    $+0$や$-0$は近づく側を表す記法であり,異なる二つの数を表しているわけではない.

    左右の両側から近づけることができる場合,通常の\bookterm{limit}{極限}$\lim_{x\to a}f(x)$が有限の値$L$になるための必要十分条件は,左右の\bookterm{one-sided}{片側極限}がともに存在し,同じ値$L$になることである.
    つまり,
    \[
      \lim_{x\to a^-}f(x)
      =\lim_{x\to a^+}f(x)=L
    \]
    である.
    例えば,
    \[
      f(x)=
      \begin{cases}
        -1 & (x<0),\\
        1 & (x\ge0)
      \end{cases}
    \]
    の場合,\bookterm{left-limit}{左側極限}は$-1$,\bookterm{right-limit}{右側極限}は$1$となる.
    左右の値が異なるので,\,$\lim_{x\to0}f(x)$は存在しない.
    また,\bookterm{one-sided}{片側極限}でも$x=a$での関数の値そのものは関係しない.

  \subsection{\texorpdfstring{ネイピア数$e$}{ネイピア数e}}\label{節：ネイピア数eの定義}
    ネイピア数$e$は,ベルヌーイの定義によって,
    \[
      e=\lim_{n\to\infty}\left(1+\frac{1}{n}\right)^n
    \]
    と定める.
    以下では,この定義と\bookterm{sequence}{数列}・関数の\bookterm{limit}{極限}を用いて,\bookterm{euler}{オイラーの公式}の導出に必要な関係を確認する.

  \subsection{三角関数の極限}\label{節：三角関数の基本極限}
    角度は弧度法で表す.
    $0<x<\dfrac{\pi}{2}$のとき,半径$1$の円の扇形と,その扇形に内接・外接する三角形の面積を比べると,

    \[
      \frac{1}{2}\sin x
      <
      \frac{1}{2}x
      <
      \frac{1}{2}\tan x
    \]
    となる.
    \begin{center}
      \BookMathLayout{\begin{tikzpicture}[
        x=2.4cm,
        y=2.4cm,
        line cap=round,
        line join=round,
        >=stealth
      ]
        % x=40度を例に,扇形と内接・外接する三角形を描く.
        \pgfmathsetmacro{\figangle}{40}
        \coordinate (O) at (0,0);
        \coordinate (A) at (1,0);
        \coordinate (P) at ({cos(\figangle)},{sin(\figangle)});
        \coordinate (Q) at (1,{tan(\figangle)});

        % 第一象限だけを描き,各図形は境界線で示す.
        \draw[->,thin] (O)--(1.38,0) node[below] {$x$};
        \draw[->,thin] (O)--(0,1.28) node[left] {$y$};
        \draw[thin] (A) arc[start angle=0,end angle=90,radius=1];
        \draw[thick] (O)--(Q)--(A)--cycle;
        \draw[thick] (O)--(P)--(A);
        \draw[thick] (A) arc[start angle=0,end angle=\figangle,radius=1];
        \draw[->,semithick] (0.24,0)
          arc[start angle=0,end angle=\figangle,radius=0.24];

        \fill (O) circle (0.018);
        \fill (A) circle (0.018);
        \fill (P) circle (0.018);
        \fill (Q) circle (0.018);
        \node[below left] at (O) {$O$};
        \node[below] at (0.52,-0.03) {$1$};
        \node[below right] at (A) {$A$};
        \node[above left] at (P) {$P$};
        \node[above right] at (Q) {$Q$};
        \node at (0.55,0.20) {$\frac12\sin x$};
        \node at (0.27,0.07) {$x$};

        \node[anchor=east,align=right,text width=2.0cm] (sector-label)
          at (-0.10,0.88) {扇形\\[-1pt]面積 $\frac12x$};
        \draw[->,thin] (sector-label.east)--(0.94,0.342);
        \node[anchor=west,align=left,text width=2.1cm] (outer-label)
          at (1.33,0.62) {外接三角形\\[-1pt]面積 $\frac12\tan x$};
        \draw[->,thin] (outer-label.west)--(0.91,0.55);
      \end{tikzpicture}}{\resizebox{\linewidth}{!}{\begin{tikzpicture}[
        x=2.4cm,
        y=2.4cm,
        line cap=round,
        line join=round,
        >=stealth
      ]
        % x=40度を例に,扇形と内接・外接する三角形を描く.
        \pgfmathsetmacro{\figangle}{40}
        \coordinate (O) at (0,0);
        \coordinate (A) at (1,0);
        \coordinate (P) at ({cos(\figangle)},{sin(\figangle)});
        \coordinate (Q) at (1,{tan(\figangle)});

        % 第一象限だけを描き,各図形は境界線で示す.
        \draw[->,thin] (O)--(1.38,0) node[below] {$x$};
        \draw[->,thin] (O)--(0,1.28) node[left] {$y$};
        \draw[thin] (A) arc[start angle=0,end angle=90,radius=1];
        \draw[thick] (O)--(Q)--(A)--cycle;
        \draw[thick] (O)--(P)--(A);
        \draw[thick] (A) arc[start angle=0,end angle=\figangle,radius=1];
        \draw[->,semithick] (0.24,0)
          arc[start angle=0,end angle=\figangle,radius=0.24];

        \fill (O) circle (0.018);
        \fill (A) circle (0.018);
        \fill (P) circle (0.018);
        \fill (Q) circle (0.018);
        \node[below left] at (O) {$O$};
        \node[below] at (0.52,-0.03) {$1$};
        \node[below right] at (A) {$A$};
        \node[above left] at (P) {$P$};
        \node[above right] at (Q) {$Q$};
        \node at (0.55,0.20) {$\frac12\sin x$};
        \node at (0.27,0.07) {$x$};

        \node[anchor=east,align=right,text width=2.0cm] (sector-label)
          at (-0.10,0.88) {扇形\\[-1pt]面積 $\frac12x$};
        \draw[->,thin] (sector-label.east)--(0.94,0.342);
        \node[anchor=west,align=left,text width=2.1cm] (outer-label)
          at (1.33,0.62) {外接三角形\\[-1pt]面積 $\frac12\tan x$};
        \draw[->,thin] (outer-label.west)--(0.91,0.55);
      \end{tikzpicture}}}
      \BookFigureCaption{扇形と内接・外接三角形の面積比較}
    \end{center}
    したがって,
    \[
      \cos x<\frac{\sin x}{x}<1
    \]
    である.
    $x\to+0$のとき$\cos x\to1$なので,はさみうちの原理より
    \[
      \lim_{x\to+0}\frac{\sin x}{x}=1
    \]
    を得る.
    また,\,$\dfrac{\sin x}{x}$は偶関数であるから,左側から近づけた場合にも同じ\bookterm{limit}{極限}となる.
    よって,
    \[
      \lim_{x\to0}\frac{\sin x}{x}=1
    \]
    である.

    ここで示した面積比較による証明は,初等的でよく用いられる説明である.
    ただし,扇形の面積公式や図形の大小関係を直観的に用いており,それらの根拠を厳密に示していないため,厳密な証明ではなく直観的な説明として位置づける.

    さらに,\,$\tan x=\dfrac{\sin x}{\cos x}$および$\lim_{x\to0}\cos x=1$より,
    \[
      \lim_{x\to0}\frac{\tan x}{x}
      =
      \lim_{x\to0}
      \left(
        \frac{\sin x}{x}\cdot\frac{1}{\cos x}
      \right)
      =1
    \]
    となる.

  \subsection{指数・対数の極限}\label{節：自然対数と指数関数の基本極限}
    以下では,\,$\log x$を自然\bookterm{logarithm}{対数}$\log_e x$の意味で用いる.
    \[
      \lim_{x\to0}\frac{\log(1+x)}{x}
    \]
    について考える.
    $x$が$0$に十分近く,\,$x\ne0$のとき,\bookterm{logarithm}{対数}の性質から
    \[
      \frac{\log(1+x)}{x}
      =\log\left((1+x)^\frac{1}{x}\right)
    \]
    である.右辺の\bookterm{limit}{極限}を求めるには,
    \[
      \lim_{x\to0}(1+x)^\frac{1}{x}=e
    \]
    を使う.これは,先に指数\bookterm{logarithm}{対数}の節で定めた$e$の\bookterm{limit}{極限}表示から確認できる.
    $x\to 0^+$では,\,$\displaystyle t=\frac{1}{x}\to\infty$と置くと
    \[
      (1+x)^\frac{1}{x}=\left(1+\frac1t\right)^t\longrightarrow e.
    \]
    $x\to 0^-$では,\,$\displaystyle t=-\frac{1}{x}\to\infty$と置くと
    \begin{align*}
      &(1+x)^\frac{1}{x}
      =\left(1-\frac1t\right)^{-t}\\
      &=\left(1+\frac1{t-1}\right)^{t-1}
       \left(1+\frac1{t-1}\right)
      \longrightarrow e.
    \end{align*}
    自然\bookterm{logarithm}{対数}は正の数上で連続であり,\,$\log e=1$なので,
\BookMathLayout{\[
      \lim_{x\to0}\frac{\log(1+x)}{x}
      =\lim_{x\to0}\log\left\{(1+x)^\frac{1}{x}\right\}
      =\log e=1
    \]}{
\[
\begin{aligned}
\lim_{x\to0}\frac{\log(1+x)}x
&=\lim_{x\to0}\log\left\{(1+x)^{1/x}\right\}\\
&=\log e=1
\end{aligned}
\]
}
    となる.
    この\bookterm{limit}{極限}を用いて,\,$\lim_{n\to\infty}\left(1+\frac{a}{n}\right)^n$を求める.
    $a=0$の場合は,
    \[
      \lim_{n\to\infty}\left(1+\frac{0}{n}\right)^n=1=e^0
    \]
    である.以下,\,$a\ne0$とする.
    $n$が十分大きければ,\,$1+\dfrac{a}{n}>0$なので,
    \[
      \lim_{n\to\infty}\left(1+\frac{a}{n}\right)^n
      =
      \left\{
        \lim_{n\to\infty}
        \left(1+\frac{a}{n}\right)^{\frac{n}{a}}
      \right\}^{a}
      =e^a
    \]
    となる.以上より,\,$a=0$を含む任意の実数$a$について,
    \[
      \lim_{n\to\infty}\left(1+\frac{a}{n}\right)^n=e^a
    \]
    が成り立つ.

% section: 極形式と複素平面

\section{複素平面と極形式}\label{節：極形式と複素平面}
  \noindent\textbf{【本編とのつながり】}\par
  複素数の積を回転と拡大・縮小として読む.\sectionref*{節：1次分数型}の図形的解釈と\sectionref*{節：虚数解型隣接3項間漸化式}で使う.\bookterm{polar}{極形式}の計算には\bookterm{trigonometric}{三角関数}を,\bookterm{euler}{オイラーの公式}の導出には\bookterm{binomial}{二項定理}と\bookterm{limit}{極限}を使う.

本節では,座標平面における回転,移動,縮小,拡大および\bookterm{de-moivre}{ド・モアブルの定理},\bookterm{euler}{オイラーの公式}について扱う.

  \subsection{複素数を平面に表す}\label{節：複素平面と極形式}

    $i^2=-1$を満たす虚数単位$i$を用いると,実数$x,y$に対して複素数は
    \[
      z=x+iy
    \]
    と表される.
    この複素数$z$を座標平面上の点$(x,y)$に対応させたものを\bookterm{complex-plane}{複素平面}という.
    横軸を実軸,縦軸を虚軸とする.

    点$(x,y)$と原点との距離は
    \[
      \rho=\sqrt{x^2+y^2}
    \]
    である.
    また,正の$x$軸から点$(x,y)$へ向かう半直線まで反時計回りに測った角を$\theta$とすると,
    \[
      x=\rho\cos\theta,\qquad
      y=\rho\sin\theta
    \]
    と表せる.
    したがって,\,$z$は
    \[
      z=\rho(\cos\theta+i\sin\theta)
    \]
    と表される.
    この表し方を複素数の\bookterm{polar}{極形式}という.

    $\rho$は$z$の絶対値であり,\,$|z|=\rho$と書く.
    $z\ne0$のとき,\,$\theta$を$z$の\bookterm{argument}{偏角}という.
    \bookterm{argument}{偏角}には$2\pi$の整数倍を加えても同じ点を表すため,一意には定まらない.
    $z=0$では$\rho=0$となり,\bookterm{argument}{偏角}は定まらない.
  \BookFigPolar

    例えば,\,$\rho=2,\,\theta=\dfrac{\pi}{3}$の点は
    \[
      (x,y)
      =
      \left(
        2\cos\frac{\pi}{3},
        2\sin\frac{\pi}{3}
      \right)
      =
      (1,\sqrt{3})
    \]
    である.
    直交座標では$(1,\sqrt{3})$,\bookterm{polar}{極形式}では
    \[
      2\left(
        \cos\frac{\pi}{3}
        +i\sin\frac{\pi}{3}
      \right)
    \]
    と表される.

  \subsection{複素数と平面の変換}\label{節：複素数による平面上の変換}

    複素数の加法と乗法は,座標平面上の点の移動や回転,拡大・縮小に対応する.

    まず,\,$b=u+iv$を固定した複素数とし,\,$z=x+iy$に$b$を加える.
    \[
      w=z+b=(x+u)+i(y+v)
    \]
    となるので,点$(x,y)$は点$(x+u,y+v)$へ移る.
    これは,すべての点を同じベクトル$(u,v)$だけ平行移動する操作である.

    次に,\,$z=\rho(\cos\theta+i\sin\theta)$に
    \[
      c=s(\cos\phi+i\sin\phi)
      \qquad(s>0)
    \]
    を掛ける.
    \bookterm{trigonometric}{三角関数}の\bookterm{addition}{加法定理}より,
    \begin{align*}
      cz
        &=s\rho
          (\cos\phi+i\sin\phi)
          (\cos\theta+i\sin\theta)\\
        &=s\rho
          \{\cos(\theta+\phi)+i\sin(\theta+\phi)\}.
    \end{align*}
    したがって,絶対値は$s$倍され,\bookterm{argument}{偏角}は$\phi$だけ増える.

    特に,\,$s=1$なら原点を中心とする角$\phi$の回転である.
    $s>1$なら原点からの距離が$s$倍される拡大,\,$0<s<1$なら原点からの距離が$s$倍される縮小となる.
  \BookFigComplexRotation

    例えば,\,$i=\cos\dfrac{\pi}{2}+i\sin\dfrac{\pi}{2}$であるから,\,$z$に$i$を掛けると原点を中心に反時計回りに$\dfrac{\pi}{2}$だけ回転する.

    回転・拡大縮小の中心を$z_0$とし,移動先の中心を$w_0$とすると,これらをまとめた変換は
    \[
      w=c(z-z_0)+w_0
    \]
    と表せる.
    まず中心からの差$z-z_0$に$c$を掛けて回転・拡大縮小し,その後$w_0$を加えて中心を移す.

  \subsection{ド・モアブルの定理}\label{節：ド・モアブルの定理}
    まず,\bookterm{polar}{極形式}どうしの積を考える.
    \bookterm{trigonometric}{三角関数}の\bookterm{addition}{加法定理}から,
    \begin{align*}
      &(\cos\alpha+i\sin\alpha)
       (\cos\beta+i\sin\beta)\\
      &\quad=
       \cos\alpha\cos\beta-\sin\alpha\sin\beta\\
      &\qquad
       +i(\sin\alpha\cos\beta+\cos\alpha\sin\beta)\\
      &\quad=
       \cos(\alpha+\beta)+i\sin(\alpha+\beta)
    \end{align*}
    となる.
    つまり,絶対値が$1$の複素数どうしを掛けると,\bookterm{argument}{偏角}が加わる.

    これを繰り返すと,任意の整数$n$について
    \[
      (\cos\theta+i\sin\theta)^n
      =
      \cos(n\theta)+i\sin(n\theta)
    \]
    が成り立つ.
    これを\bookterm{de-moivre}{ド・モアブルの定理}という.

    正の整数$n$については,\bookterm{induction}{数学的帰納法}で確かめられる.
    $n=1$では明らかであり,\,$n$で成り立つと仮定すると,
    \begin{align*}
      &(\cos\theta+i\sin\theta)^{n+1}\\
        &\qquad =\{\cos(n\theta)+i\sin(n\theta)\}
        (\cos\theta+i\sin\theta)\\
        &\qquad =\cos((n+1)\theta)+i\sin((n+1)\theta)
    \end{align*}
    となる.
    また,\,$n=0$では両辺とも$1$である.
    $n<0$の場合は,\bookterm{addition}{加法定理}より
    \[
      (\cos\theta+i\sin\theta)
      (\cos(-\theta)+i\sin(-\theta))=1
    \]
    であるから,
    \[
      (\cos\theta+i\sin\theta)^{-1}
      =
      \cos(-\theta)+i\sin(-\theta)
    \]
    を用いて正の整数の場合に帰着できる.

    例えば,
    \[
      \left(
        \cos\frac{\pi}{6}
        +i\sin\frac{\pi}{6}
      \right)^6
      =
      \cos\pi+i\sin\pi
      =
      -1
    \]
    である.
    このように\bookterm{polar}{極形式}で表すと,累乗では絶対値を累乗し,\bookterm{argument}{偏角}を指数倍すればよい.

  \subsection{オイラーの公式}\label{節：オイラーの公式}
    次に,複素数の指数表示$e^{i\theta}$を\bookterm{limit}{極限}を用いて考える.
    ここでは
    \[
      e^{i\theta}
      =
      \lim_{n\to\infty}
      \left(1+\frac{i\theta}{n}\right)^n
    \]
    と定める.
    前節で準備した\bookterm{limit}{極限}を用いて,この値を求める.

    $\theta\ne0$とし,\,$\phi_n$を
    \[
      -\frac{\pi}{2}<\phi_n<\frac{\pi}{2},
      \qquad
      \tan\phi_n=\frac{\theta}{n}
    \]
    を満たす角とする.
    また,
    \[
      r_n=\sqrt{1+\frac{\theta^2}{n^2}}
    \]
    とおくと,
    \[
      \cos\phi_n=\frac{1}{r_n},
      \qquad
      \sin\phi_n=\frac{\theta}{nr_n}
    \]
    である.
    よって,
    \[
      1+\frac{i\theta}{n}
      =
      r_n(\cos\phi_n+i\sin\phi_n)
    \]
    と\bookterm{polar}{極形式}で表せる.

    まず,\,$r_n^n$の\bookterm{limit}{極限}を求める.
    自然\bookterm{logarithm}{対数}の基本\bookterm{limit}{極限}より,
    \begin{align*}
      \log(r_n^n)
        &=
        \frac{n}{2}
        \log\left(1+\frac{\theta^2}{n^2}\right)\\
        &=
        \cfrac{\theta^2}{2n}
        \cfrac{
          \log\left(1+\cfrac{\theta^2}{n^2}\right)
        }{
          \cfrac{\theta^2}{n^2}
        }\\
        &\longrightarrow0.
    \end{align*}
    したがって,
    \[
      r_n^n\longrightarrow1
    \]
    である.

    次に,\,$n\phi_n$の\bookterm{limit}{極限}を求める.
    $\tan\phi_n=\dfrac{\theta}{n}\to0$より$\phi_n\to0$であり,
    \[
      n\phi_n
      =
      \theta\frac{\phi_n}{\tan\phi_n}
      \longrightarrow\theta
    \]
    となる.
    ここでは,\,$\lim_{x\to0}\dfrac{x}{\tan x}=1$を用いた.

    \bookterm{de-moivre}{ド・モアブルの定理}により,
    \[
      \left(1+\frac{i\theta}{n}\right)^n
      =
      r_n^n
      \{\cos(n\phi_n)+i\sin(n\phi_n)\}
    \]
    である.
    $r_n^n\to1,\,n\phi_n\to\theta$と\bookterm{trigonometric}{三角関数}の連続性から,
    \[
      e^{i\theta}
      =
      \lim_{n\to\infty}
      \left(1+\frac{i\theta}{n}\right)^n
      =
      \cos\theta+i\sin\theta
    \]
    を得る.
    $\theta=0$の場合も両辺が$1$であり,同じ式が成り立つ.
    これを\bookterm{euler}{オイラーの公式}という.

    \bookterm{euler}{オイラーの公式}により,複素数の\bookterm{polar}{極形式}
    \[
      z=\rho(\cos\theta+i\sin\theta)
    \]
    は
    \[
      z=\rho e^{i\theta}
    \]
    とも表せる.
    $z\ne0$,すなわち$\rho>0$なら,負の整数乗も逆数を用いて定められる.
    この条件の下で,\bookterm{de-moivre}{ド・モアブルの定理}は
\BookMathLayout{\[
      z^n
      =
      \rho^n e^{in\theta}
      =
      \rho^n\{\cos(n\theta)+i\sin(n\theta)\}
      \qquad(n\in\mathbb Z)
    \]}{
\[
\begin{aligned}
z^n&=\rho^n e^{in\theta}\\
&=\rho^n\{\cos(n\theta)+i\sin(n\theta)\}\\
&\hspace{5zw}(n\in\mathbb Z)
\end{aligned}
\]
}
    と書ける.\,$z=0$では負の整数乗を定めず,正の整数$n$について$0^n=0$とする.

  \subsection{複素数の実数乗}\label{節：複素数の実数乗}
    $z=\rho e^{i\theta}$で$\rho>0$とし,\bookterm{argument}{偏角}$\theta$を一つ選んで固定する.
    この選び方に基づいて,実数$x$に対する複素数のべき乗を
    \[
      z^x=\rho^x e^{ix\theta}
      =
      \rho^x\{\cos(x\theta)+i\sin(x\theta)\}
    \]
    と定めることができる.

    ただし,\bookterm{argument}{偏角}は$\theta+2\pi k$とも表せるため,整数でない$x$に対する$z^x$の値は,どの\bookterm{argument}{偏角}を選ぶかによって変わる.
    $x$が整数なら,選ぶ\bookterm{argument}{偏角}によらず同じ値になる.
    また,\,$x=\dfrac{m}{n}$を有理数とし,\,$n\ge1$として上の定義で得た値を$n$乗すると,
    \[
      (z^{\frac{m}{n}})^n=z^m
    \]
    となる.
    したがって,この定義で得られる$z^{\displaystyle \frac{m}{n}}$は$z^m$の$n$乗根の一つである.

    \subsection{公式のまとめ}\label{節：議論のまとめ}
    ここまでの議論が長くなったので,扱った内容をまとめる.

    複素数$z=x+iy$を座標平面上の点$(x,y)$に対応させると,複素数を平面上の点として考えられる.
    \bookterm{polar}{極形式}
    \[
      z=\rho(\cos\theta+i\sin\theta)
    \]
    では,\,$\rho$が原点からの距離を表し,\,$\theta$が\bookterm{argument}{偏角}を表す.

    複素数に一定の数を加えることは平行移動に対応し,\bookterm{polar}{極形式}で表した複素数を掛けることは回転と拡大縮小に対応する.
    積を繰り返すと\bookterm{argument}{偏角}が加わるため,\bookterm{de-moivre}{ド・モアブルの定理}が得られる.
    また,\,\sectionref{節：極限}で準備した\bookterm{limit}{極限}を用いると\bookterm{euler}{オイラーの公式}
    \[
      e^{i\theta}=\cos\theta+i\sin\theta
    \]
    が導かれる.
    このように,複素数を\bookterm{polar}{極形式}で表すことで,座標平面上の変換と複素数の計算を結び付けられる.

\end{bookcolumns}
